\documentclass[12pt]{article}
\usepackage[utf8]{inputenc}
\usepackage[T1]{fontenc}
\usepackage[protrusion=true,expansion=true]{microtype}
\usepackage[margin=1.4in]{geometry}
\usepackage{setspace}
\usepackage{libertinus}
\usepackage{libertinust1math}
\usepackage[colorlinks=true,linkcolor=black,urlcolor=black,
  pdftitle={On Immortality},
  pdfauthor={Soren Kierkegaard}]{hyperref}
\setlength{\parindent}{1.5em}
\setlength{\parskip}{0pt}
\onehalfspacing

\title{On Immortality\\[0.5em]
  \large\textit{Concluding Unscientific Postscript}, pp.\ 153--158}
\author{S\o{}ren Kierkegaard\\[0.4em]
  \small Translated by David~F.\ Swenson and Walter Lowrie}
\date{}

\begin{document}
\maketitle

\noindent
When death thus becomes something to be related to the entire life of
the subject, I must confess I am very far indeed from having understood
it, even if it were to cost me my life to make this confession. Still
less have I realized the task existentially. And yet I have thought
about this subject again and again; I have sought for guidance in
books---and I have found none.\footnote{Although I have said this
often, I wish here again to repeat it: what is developed in these
pages does not concern the simple-minded, who bear feelingly the
burdens of life, and whom God wishes to preserve in their lovable
simplicity, which feels no great need of any other sort of
understanding. Or in so far as such need is felt, it tends to reduce
itself to a sigh over the ills of life, the sigh humbly finding solace
in the thought that the real happiness of life does not consist in
having knowledge. On the other hand it does concern those who deem
themselves possessed of leisure and talent for a deeper inquiry. And
it concerns such an one in the following manner: it seeks to estop him
from thoughtlessly taking on universal rotatory motion, without first
considering in self-reflection that being an existing human individual
is so strenuous and yet so natural a task for everyone, that one tends
first as a matter of courage to apply himself to this task, and
reasonably finds in the exertion thereto requisite, a sufficiency for
his entire life.}

For example, what does it mean to be immortal? In this respect, I know
what people generally know. I know that some hold a belief in
immortality, that others say they do not hold it; whether they actually
do not hold it I know not; it does not occur to me therefore to want
to combat them, for such an undertaking is so dialectically difficult
that I should need a year and a day before it could become dialectically
clear to me whether there is any reality in such a contest; whether the
dialectic of communication, when it is properly understood, would
approve of such a proceeding or transform it into a mere beating of the
air; whether the consciousness of immortality is a doctrinal topic
which is appropriate as a subject for instruction, and how the dialectic
of instruction must be determined with relation to the learner's
presuppositions; whether these presuppositions are not so essential that
the instruction becomes a deception in case one is not at once aware of
them, and in that event the instruction is transformed into
non-instruction. Moreover, I know that some have found immortality in
Hegel, others have not; I know that I have not found it in the System,
where indeed it is also unreasonable to seek it; for, in a fantastic
sense, all systematic thinking is \textit{sub specie aeterni}, and to
that extent immortality is there in the sense of eternity, but this
immortality is not at all the one about which the question is asked,
since the question is about the immortality of a mortal, which is not
answered by showing that the eternal is immortal, and the immortality
of the eternal is a tautology and a misuse of words. I have read
Professor Heiberg's \textit{Soul after Death}, indeed I have read it
with the commentary of Dean Tryde. Would I had not done so; for one
rejoices aesthetically in a poetic word and does not require the utmost
dialectical precision which is appropriate in the case of a learner who
would direct his life in accordance with such guidance. If a commentator
compels me to seek for something of that sort in the poem, he has done
no service to the poem. From the commentator I could perhaps hope to
learn what I have not learnt from the commentary, if Dean Tryde with
catechetical instruction were to take pity on me and show how one
constructs a life-view upon the profundities he has propounded by
paraphrasing the poem. For honor be to Dean Tryde---merely from that
little article it would surely be possible to construct several
different life-views---but I cannot make one out of it, and that, alas,
is precisely the misfortune, for it is one I need, no more, since I am
not a learned man. I know, moreover, that the late Professor Poul
M\o{}ller, who was well acquainted after all with the newest philosophy,
became aware only in his latest period of the infinite difficulty of
the question of immortality, when the question is simply put, when one
does not ask about a new proof and about the opinions of Tom, Dick and
Harry strung upon a thread. I know also that in a treatise he sought
to explain himself, and that this treatise clearly bears the mark of
his aversion to speculation. The difficulty of the question arises
precisely when it is simply put, not as the well-trained
\textit{Privatdocent} enquires about the immortality of man, the
abstractly understood man in general, man being understood fantastically
as the race, and so about the immortality of the human race. Such a
well-trained \textit{Privatdocent} raises and answers the question in
such a way as the well-trained reader conceives that it ought to be
done. A poor untrained reader is only made a fool of by such
reflections, like one who overhears an examination in which the
questions and answers have been agreed upon beforehand, or like one who
enters a family circle which has its own language, using words of the
mother tongue, but understanding something different by them. A book
raises the question of the immortality of the soul. The contents of the
book constitute the answer. But the contents of the book, as the reader
can convince himself by reading it through, are the opinions of the
wisest and best men about immortality, all neatly strung on a thread.
Oh! thou great Chinese god! Is this immortality? So then the question
about immortality is a learned question. All honor to learning! All
honor to him who can handle learnedly the learned question of
immortality! But the question of immortality is essentially not a
learned question, rather it is a question of inwardness, which the
subject by becoming subjective must put to himself. Objectively the
question cannot be answered, because objectively it cannot be put,
since immortality precisely is the potentiation and highest development
of the developed subjectivity. Only by really willing to become
subjective can the question properly emerge, therefore how could it be
answered objectively? The question cannot be answered in social terms,
for in social terms it cannot be expressed, inasmuch as only the
subject who wills to become subjective can conceive the question and
ask rightly, ``Do I become immortal, or am I immortal?'' Of course,
people can combine for many things; thus several families can combine
for a box at the theater, and three single gentlemen can combine for a
riding horse, so that each of them rides every third day. But it is not
so with immortality; the consciousness of my immortality belongs to me
alone, precisely at the moment when I am conscious of my immortality I
am absolutely subjective, and I cannot become immortal in partnership
with three single gentlemen in turn. People who go about with a paper
soliciting the endorsement of numerous men and women, who feel a need
in general to become immortal, get no reward for their pains, for
immortality is not a possession which can be extorted by a list of
endorsements. Systematically, immortality cannot be proved at all. The
fault does not lie in the proofs, but in the fact that people will not
understand that viewed systematically the whole question is nonsense,
so that instead of seeking outward proofs, one had better seek to
become a little subjective. Immortality is the most passionate interest
of subjectivity; precisely in the interest lies the proof. When for the
sake of objectivity (quite consistently from the systematic point of
view), one systematically ignores the interest, God only knows in this
case what immortality is, or even what is the sense of wishing to
prove it, or how one could get into one's head the fixed idea of
bothering about it. If one were systematically to hang up immortality
on the wall, like Gessler's hat, before which we take off our hats as
we pass by, that would not be equivalent to being immortal or to being
conscious of one's immortality. The incredible pains the System takes
to prove immortality is labor lost and a ludicrous contradiction---to
want to answer systematically a question which possesses the remarkable
trait that systematically it cannot be put. This is like wanting to
paint Mars in the armor which rendered him invisible. The very point
lies in the invisibility; and in the case of immortality, the point
lies in the subjectivity and in the subjective development of the
subjectivity.

Quite simply therefore the existing subject asks, not about immortality
in general, for such a phantom has no existence, but about his
immortality, about what it means to become immortal, whether he is
able to contribute anything to the accomplishment of this end, or
whether he becomes immortal as a matter of course, or whether he is
that and can become it. In the first case, he asks what significance
it may have, if any, that he has let time pass unutilized, whether
there is perhaps a greater and a lesser immortality. In the second
case, he asks what significance it may have for the whole of his human
existence that the highest thing in life becomes something like a
prank, so that the passion of freedom within him is relegated to lower
tasks but has nothing to do with the highest, not even negatively, for
a negative employment with relation to the highest thing would in a
way be the most extenuating---when one wanted to do everything
enthusiastically with all one's might, and then to ascertain that the
utmost one can do is to maintain a receptive attitude towards that
thing which one would more than gladly do everything to earn. The
question is raised, how he is to comport himself in talking about his
immortality, how he can at one and the same time talk from the
standpoint of infinity and of finiteness and think these two together
in one single instant, so that he does not say now the one and now the
other; how language and all modes of communication are related thereto,
when all depends upon being consistent in every word, lest the little
heedless supplementary word, the chatty subordinate phrase, might
intervene and mock the whole thing; where may be the place, so to
speak, for talking about immortality, where such a place exists, since
he well knows how many pulpits there are in Copenhagen, and that there
are two chairs of philosophy, but where the place is which is the unity
of infinitude and finiteness, where he, who is at one and the same time
infinite and finite, can talk in one breath of his infinitude and his
finiteness, whether it is possible to find a place so dialectically
difficult, which nevertheless is so necessary to find. The question is
raised, how he, while he exists, can hold fast his consciousness of
immortality, lest the metaphysical conception of immortality proceed to
confuse the ethical and reduce it to an illusion; for ethically,
everything culminates in immortality, without which the ethical is
merely use and wont, and metaphysically, immortality swallows up
existence, yea, the seventy years of existence, as a thing of naught,
and yet ethically this naught must be of infinite importance. The
question is raised, how immortality practically transforms his life; in
what sense he must have the consciousness of it always present to him,
or whether perhaps it is enough to think this thought once for all,
whether it is not true that, if the answer is to this effect, the
answer shows that the problem has not been stated, inasmuch as to such
a consciousness of immortality once for all there would correspond the
notion of being a subject as it were in general, whereby the question
about immortality is made fantastically ludicrous, just as the converse
is ludicrous, when people who have fantastically made a mess of
everything and have been every possible sort of thing, one day ask the
clergyman with deep concern whether in the beyond they will then really
be the same---after never having been able in their lifetime to be the
same for a fortnight, and hence have undergone all sorts of
transformations. Thus immortality would be indeed an extraordinary
metamorphosis if it could transform such an inhuman centipede into an
eternal identity with itself, which this ``being the same'' amounts to.
He asks about whether it is now definitely determined that he is
immortal, about what this determinateness of immortality is; whether
this determinateness, when he lets it pass for something once for all
determined (employing his life to attend to his fields, to take to
himself a wife, to arrange world-history) is not precisely
indeterminateness, so that in spite of all determinateness he has not
got any further, because the problem is not even conceived, but since
he has not employed his life to become subjective, his subjectivity has
become some sort of an indeterminate something in general, and that
abstract determinateness has become therefore precisely
indeterminateness; whether this determinateness (if he employs his
life to become subjective) is not rendered so dialectically difficult,
by the constant effort to adapt himself to the alternation which is
characteristic of existence, that it becomes indeterminateness;
whether, if this is the highest he attains to (namely, that the
determinateness becomes indeterminateness), it were not better to give
the whole thing up; or whether he is to fix his whole passion upon the
indeterminateness, and with infinite passionateness embrace the
indeterminateness of the determinate; and whether this might be the
only way by which he can attain knowledge of his immortality so long
as he is existing, because as exister he is marvelously compounded, so
that the determinateness of immortality can only be possessed
determinately by the Eternal, but by an exister can be possessed only
in indeterminateness.

And the fact of asking about his immortality is at the same time for
the existing subject who raises the question a deed---as it is not, to
be sure, for absent-minded people who once in a while ask about the
matter of being immortal quite in general, as if immortality were
something one has once in a while, and the question were some sort of
thing in general. So he asks how he is to behave in order to express
in existence his immortality, whether he is really expressing it; and
for the time being, he is satisfied with this task, which surely must
be enough to last a man a lifetime since it is to last for an eternity.
And then? Well, then, when he has completed this task, then comes the
turn for world-history. In these days, to be sure, it is just the
other way round: now people apply themselves first to world-history,
and therefore there comes out of this the ludicrous result (as another
author has remarked) that while people are proving and proving
immortality quite in general, faith in immortality is more and more
diminishing.

\end{document}
